\section{Statement, mechanism, and status}
Fix a modulus \(q\), divisible by \(18480\), and a prime
\(d\equiv q/2-1\pmod q\). Put \(\kappa=(\,\cdot\,/d)\), the quadratic
character of conductor \(d\). The analytic goal is the bound
\[
 R_{q,d}(x)\le C_{q,d,\varepsilon}x^{-1/2+\varepsilon}
 \qquad (x\text{ sufficiently large}),
 \tag{1}
\]
for the corrected prime sum defined in the next section. The modulus
and auxiliary prime are fixed before \(x\) grows.

The proof is organized around one arithmetic comparison. Divisibility
produces short integral theta images. A moment estimate selects almost
a basis from their translates. Separate twists at the coefficient
primes bound the determinant of that same family from the opposite
direction. A common term cancels and leaves (1). The number of
successive theta lifts is of order \(\sqrt{x}\).

This is a reconstruction of the argument of Levent Alp\"oge with Claude,
using the \href{\sourcepin}{Sernau--Anthropic rewriting} as the principal
source. Our mathematical proofs are distinct from the published formal
verification. We do not claim that the present ordinary reconstruction
is complete. In particular the arithmetic constructions in Sections 3--5
are open as indicated. The detailed working proofs, not Lean declarations
or assertions about source length, are the \emph{mathematical} references.

\section{An exact analytic criterion}
For each character \(\psi\bmod q\), let \(E_\psi(s)\) be \(L(s,\psi)\)
with its pole removed by multiplying by \(s-1\) in the principal case.
Set
\[
 P(s)=\prod_{\psi\bmod q}E_\psi(s),\qquad
 Q(s)=\prod_{\psi\bmod q}L(s,\psi\kappa),\qquad
 G(s)=\frac{Q'/Q(s)-P'/P(s)-1}{\varphi(q)}.
\]
These products have no pole on the right half-plane. Put \(c=G(1)\)
and
\[
 R_{q,d}(x)=\sum_{\substack{n\le x\\n\equiv1\pmod q}}
 \Lambda(n)(1-\kappa(n))(1/n-1/x)
 -\frac{\log x}{\varphi(q)}-c.
\]
Orthogonality and termwise integration in \(\Rea s>1\) give
\[
 \int_1^\infty R_{q,d}(x)x^{-s}\,dx
 =\frac{G(s)/s-c}{s-1}.
 \tag{2}
\]
Indeed the integral of \((1/n-1/x)x^{-s}\) from \(n\) to infinity
is \(n^{-s}/[s(s-1)]\). The logarithmic derivative has principal
term \(s/[\varphi(q)(s-1)]\) after subtracting \(G\).
The singularity of (2) at 1 is removable. At a zero \(\rho\ne1\)
in \(\Rea s>0\) the residue is
\[
 \frac{\ord_\rho P-\ord_\rho Q}{\varphi(q)\rho(1-\rho)}.
 \tag{3}
\]

\begin{lemma}[One-sided zero detection]
Let \(1/2<\beta<1\). Suppose \(R_{q,d}(x)=O_{\rm upper}
(x^{\beta-1})\) and \(\ord_t Q\le\ord_t P\) for real
\(\beta<t<1\). Then the products \(P,Q\) have identical zero
multiplicities in \(\Rea s>\beta\).
\end{lemma}
\begin{proof}
Writing \(x=e^u\), subtract the transform in (2) from
\(K/(s-\beta)\) for a sufficiently large \(K\). The difference is
the Laplace transform of the nonnegative density
\(Ke^{\beta u}-e^uR_{q,d}(e^u)\). It converges for \(\Rea s\)
large. If its real abscissa of convergence \(a\) exceeded \(\beta\),
positivity would make the residue at \(a\) nonnegative, whereas
(3) makes it nonpositive. It is zero, so \(a\) would be a regular
point. This is impossible for a positive Laplace transform:
the Taylor series at a real point slightly to the right of \(a\),
expanded towards the left, has nonnegative integral terms.
Convergence past \(a\) would imply convergence of the defining
integral there by monotone convergence. Thus \(a\le\beta\).
Equation (2) is holomorphic to the right of \(\beta\);
(3) proves the assertion.
\end{proof}

The \prooflink{notes/lsernau-zero-detection.md}{finite-box
quadratic-twist selection theorem} is proved in ordinary mathematics:
for fixed \(q,\beta,H\), all but
\(O_{q,\beta,H}(X^{1-(\beta-1/2)/18})\) primes \(d\le X\)
in the chosen progression give no zeros of any \(L(s,\psi\kappa)\)
on \(\Rea s\ge\beta,\ |\Im s|\le H\).
Its proof uses the character theta transformation, a truncated
M\"obius Mellin integral, three-lines, and the Jacobi second
moment from quadratic reciprocity and P\'olya--Vinogradov. The
reciprocal sum over the bad primes converges; the reciprocal sum
over the progression diverges. Hence a good prime exists.

For a proposed zero of \(P\) to the right of \(1/2\), choose
\(\beta\) below its real part and then a good \(d\). Bound (1)
with \(\varepsilon<\beta-1/2\) contradicts the lemma.
The functional equation excludes zeros on the other side.
The imprimitive Euler factors are nonzero in \(\Rea s>0\).
Thus (1) for the admissible fixed moduli and primes suffices.
The proof of selection and the exact constant in (2) belong to
the cited mathematical chapter, not to the missing arithmetic step.

\section{The common arithmetic data}
Work on one admissible set of the readable source parameters:
\(E=\Q(\sqrt{-d},\zeta_{q_1})\), \(m\) even,
\(k=r+1-m\), \(t=5r\). The upper definite spaces have ranks
\(n_j=k+j(t+1)\); the lower Bessel space has rank \(r\), and
the lower chain has ranks \(r+jt\), ending in the definite
space of rank \(n_m-1\). The
\prooflink{notes/lsernau-common-theta-finite.md}{common finite theta chapter}
specifies the signs, parameters, principal levels at the two
chosen split primes, and the last-row level at the auxiliary prime
\(N\). On the upper chain these levels have trivial stabilizers:
a finite stabilizer embeds into an odd-prime principal congruence
subgroup of a general linear group, which is torsion-free.
Consequently the jointly invariant combination of upper theta
kernels is an exact finite sum of products of jet maps over
tuples of class representatives, with no stabilizer denominators.
This follows by successive class integration, without requiring
individual summands of the tensor to be level-invariant.

Use the Riesz representative \(\mathcal P\) of the Bessel period
and the scaled map \(P=(\mathrm{vol}J\,\nu)^{-1/2}\mathcal P\)
in the prescribed source coordinates. Let \(E_0\) denote the
central good-type projector and \(\rho(s)\) the finite level action.
At a coefficient embedding \(\iota\), the target lattice consists
of vectors whose lower-chain images \(\Theta_\xi y\) are integral
for every allowed twist. The normalized maps used below are
\[
 y_s=c_Bc_{\rm ss}^{-1}|\mathcal G|\rho(s)E_0\mathcal P\phi,
 \qquad
 c_B=\prod_{\ell\in S}\ell^{-\lfloor g_\ell-e_\ell\rfloor}.
\]
Here \(S\) is the source's set of gain and error primes,
including \(N\), \(c_{\rm ss}\ne0\) is the common see-saw
scalar, and \(g_\ell,e_\ell\) are the required local gains
and losses. The finite consequences of this normalization
are proved in the linked chapter.

The analytic comparison has three separate steps. First, the
\prooflink{notes/lsernau-global-cut-unfolding.md}{class-return identity}
turns the Bessel orbit through \(xJ\) into a finite cut with multiplicity
\([S_x:S_x\cap S_0]\), where \(S_x=R_f\cap xJx^{-1}\) and
\(S_0=R_f\cap K_f\). Its measure is
\(\mathrm{vol}(J)/\mathrm{vol}(S_0)\); the transpose index cancels the
twist-dependent factor. Second, the
\prooflink{notes/lsernau-graded-period-unfolding.md}{graded unfolding}
proves absolute convergence against rapid inputs, the surviving-orbit
classification and the one-step period identity. Third,
\prooflink{notes/lsernau-terminal-period-matching.md}{terminal matching}
identifies the final branching pairing and common scalar. Its raw scalar
may vanish, so these identities alone do not supply nonzero chain volumes.
The \prooflink{notes/lsernau-support-to-cut.md}{support theorem} supplies
the finite cuts and their prime-to-twist integral weights.

The \prooflink{notes/lsernau-bessel-majorants.md}{Bessel majorants} and
\prooflink{notes/lsernau-stationary-bessel.md}{stationary period}
supply the normalized real coefficient, its positive lower period and
the cut-tail estimate, with normalized majorant exponent \(119r^2\).
The \prooflink{notes/lsernau-rallis-sharp-inversion.md}{convergent Rallis
identity} matches the doubled Schwartz normalization, bounds the finite
vacuum inverse and gives injectivity for a supplied nonzero real edge.
The full see-saw, compatible nonzero volumes and admissible-scale
existence are retained inputs.

\paragraph{The noncompact local input.}
For the actual shifted parameters, the
\prooflink{notes/lsernau-noncompact-real-edges.md}{highest polynomial}
is a product of nested minors. Fischer pairing proves uniqueness and
polynomial saturation. The
\prooflink{notes/lsernau-discrete-series-suppliers.md}{classical coefficient
and rigidity theorem} supplies its positive polar formula and identifies
the discrete series: the bottom PRV component and Dirac inequality give
\(c\le\|\lambda\|^2-\|\rho\|^2\), while the required reverse
Casimir inequality forces equality and Vogan--Zuckerman rigidity.
The \prooflink{notes/lsernau-real-grid-parameters.md}{parameter check}
verifies regularity, both determinant shifts and the required gaps.

For the projected harmonic \(x=P_l\Phi\), row rank \(r\), column
rank \(m\), degree \(D\), and \(|l_i|\le L\), the
\prooflink{notes/lsernau-projection-lower-bound.md}{quantitative bound} is
\[
 \frac{\|P_l\Phi\|_B^2}{\|\Phi\|_F^2}
 \ge\frac{(48T)^{-2r^2}}{4\dim\tau_\Lambda},
 \qquad T=D+rm+2r(L+r+1)+1.
\]
Positive polar averaging and coefficient decay cancel the Haar scale.
The \prooflink{notes/lsernau-quantitative-local-smearing.md}{smearing
construction} gives, on each finite polynomial-degree range, a linear
map \(\mathcal S_B\) with common compact support such that
\(\pi'(\mathcal S_BP)x=P_lP\). Integral spanning matrices and
Cramer's rule give explicit differential order, denominators and smooth
seminorms in the original Fischer norm. This controls the harmonic line
and finite polynomial input; the inverse on the entire projected
quotient and the global cost of these noncompact operations remain open.
They are separate from the definite chain transport below.

\section{Short forms and enough forms}

\paragraph{One local theta package.}
The \prooflink{notes/lsernau-cm-theta-normal-form.md}{CM normal form},
\prooflink{notes/lsernau-weil-normalization.md}{standard normalization}
and \prooflink{notes/lsernau-charged-jet-estimate.md}{charged jet theorem}
give integral scalar sections at every prime, the actual trace-self-dual
lattice and vacuum, and exact Fourier and product laws. At 2 the
complementary CM idempotents keep the marked point nonzero.
The \prooflink{notes/lsernau-nonstandard-weil-transport.md}{Weil transport}
fixes the finite and real phases, proves rational unitary invariance and
canonical orthogonal-product transport, and supplies absolute mixed
summation. The
\prooflink{notes/lsernau-common-jet-conjugation.md}{common jet comparison}
then transports these data through every complex embedding; its finite
multiplier has modulus one. The definite Gaussian estimate is
\[
 |\Theta(\rho(y)\omega(g)v)|
 \le e^{a_0/8}(8/a_0)^{nf/2}\sqrt{(nf)!},
 \qquad a_0=\pi/\sqrt d.
\]
The Fock integral gives the conjugated Schur-jet norm. For a filling of
degree \(B\), charges \(m_0\), and scalar parameter \(\varpi\), the
corresponding valuation estimate is
\[
 v(\operatorname{jet})\ge
 (1/(p-1)-v(\varpi))B-m_0-v(\pi_0)D^*
                    -S_{\max}/(p-1).
\]
Here \(D^*\) is a finite minimum-cut cost and \(S_{\max}\) is the
exact maximal digit sum on the differential support.

The \prooflink{notes/lsernau-marked-chain-charges.md}{charge comparison}
cancels the diagram weights on conjugate fibres. Moving each local
Hecke polynomial through the whole prefix before taking coefficients
saves half its transfer allowance. The
\prooflink{notes/lsernau-split-transfer-and-level-lifts.md}{split transfer}
proves the exact local polynomial, auxiliary gluing and Up identities,
with last-column Fourier cost \(N^{fn_{j+1}}\).
The \prooflink{notes/lsernau-exceptional-polynomial-geometry.md}{outer}
and \prooflink{notes/lsernau-exceptional-middle-edge.md}{middle}
edge estimates control the integral elementary divisors; the
\prooflink{notes/lsernau-marked-hecke-index.md}{marked index theorem}
propagates them through the repeated product. The resulting truncated
colength is at most \(\dim X(6fa/p+1/m)\), replacing
\(\dim X(14fa/p+1/m)\); ramified intersection adds no cost.
The \prooflink{notes/lsernau-branching-normalization.md}{terminal
branching matrix} is integral, with norm at most \(L^{3L}\) when
\(L\) bounds its affected boxes. These are local inputs to one
chain calculation, rather than separate global integrality arguments.

\paragraph{Transport once through the genuine chain.}
The \prooflink{notes/lsernau-common-chain-transport.md}{common-chain
theorem} uses the actual adjacent hom spaces, not a tensor power of the
scalar sharpness test. Write \(\delta\bar\delta=M\). The standard
datum has denominator \(\delta\), while the jet scale is
\(\bar\delta j_0\); these cancel in every transported finite gauge.
An edge of degree \(h_j\) therefore loses only
\(h_jv_p(d)/2\), including at primes dividing \(M\).

The auxiliary gluing has exactly
\[
 D_N=N^{f\sum_{i<m/2}n_{2i+1}}
     =N^{fm(n_0+n_m)/4}
\]
distinct translated terms. Their supports above \(N\) separate them
from the original translations, so their coefficients and phases
are unchanged. Vacuum compression is a left inverse of gluing; on
the vacuum sector the gluing Gram operator is \(D_N I\).
Joint tensor invariance suffices for the finite class-sum formula.
Consequently the genuine composite is integral away from the twist
and \(Nd\), and its nonsplit charged gain is the sum of the already
proved edge gains, with the stated collision correction.

There is also one explicit norm bound. Let \(H_j\) be the actual
class counts, \(C_j\) the proved uniform conjugated edge norms,
\(C_\sigma=\sum_x|\sigma\Psi(x)|\), and \(c_{\rm br}\) the
integral branching norm. If the initial squared norm is at most
\(B_0>0\), then
\[
 \|\sigma\operatorname{compositeLift}(1,\Psi,\phi)(y)\|^2
 \le B_0c_{\rm br}^2N^{-b_*}D_N^2C_\sigma^2
                          \prod_{j<m}H_j^2C_j^2.
\]
Here \(b_*\) is the sum of the terminal branching exponents.
Multiplication of the edge matrices and the triangle inequality
prove the formula directly. The
\prooflink{notes/lsernau-common-chain-norm-budget.md}{norm-budget theorem}
evaluates every term on the actual scales. Its short class-count
argument records a bounded integral Gram matrix and an intermediate
module \(O_E^n\subset L\subset G^{-1}O_E^n\); Minkowski's second
theorem and finite subgroup counting give \(\log H_j\le x^{60}\).
All costs outside the row weights total at most \(fa\log b\).
The exact diagram sum and the prescribed level window then give
\[
 \log\frac{\|\sigma\operatorname{compositeLift}(1,\Psi,\phi)(y)\|^2}
                 {B_0}
 \le H_{\rm main}-mR+2\sum_pG_p\log p+108fa\log b.
\]
For zero output use the equivalent exponential inequality.
This is stronger than the required error \(fa(\log b)^3\).
It removes both the separate asymptotic norm obligation and the
quantitative adelic approximation previously used to count classes.
The analytic see-saw datum and exceptional-prime corrections remain.

\paragraph{Sharp scalar tests at every prime.}
The \prooflink{notes/lsernau-sharp-scalar-couplings.md}{scalar sharpness
theorem} supplies the least-valuation bound on the actual admissible
weights and the fixed CM period scale. Put
\(n_\sharp=\lceil6b/k\rceil\) and
\(B_\tau=\prod_{j=1}^{n_\sharp}((\ell_\tau-n_\sharp)_++j)\).
There are six scalar integral theta sections, of total multidegree
\(\ell\), with product coefficient \(a_\ell\) satisfying
\[
 v(\ell!)+v(a_\ell)\le S_p+2\sum_\tau v(B_\tau),\qquad v(p)=1.
\]
The formal-group derivative minima are superadditive. Their values
at powers of the residue cardinality give the sharp scalar bound.
An isotropic sum of three squares at odd primes, or five squares
at 2, cancels the theta factor integrally. At inert primes, a
logarithmic-derivative argument extracts the unique coupling term;
sparse near pairs have disjoint ends, so their factorial cost is
one binomial carry count per coordinate. The proof checks the
small- and large-prime weight cases, the coprime frame's unit row,
its multiplicities and its size against the stated admissible
inequalities.
The \prooflink{notes/lsernau-six-copy-test-transport.md}{actual
four-block comparison} now uses these six copies directly. Keep
the same target \(U(2z)\), and replace block sizes
\((12,4,4,2z-20)\) by \((6,4,4,2z-14)\).
The integral monomial vector has the exact multinomial Gram factor,
block theta multiplication preserves the canonical phase, and the
three auxiliary flat class sums have the same stabilizer and volume
normalizations. Their total vacuum exponent remains \(2z\).
The sharp factor is therefore genuinely
\(t_6=M^{6f}c_6\prod_\tau B_\tau^2\), with half the previous
Fourier exponent. Its actual rank-one false-lift orthogonality is
proved in the improved range \(33<2z\), as follows.

The \prooflink{notes/lsernau-rank-one-false-lift-periods.md}{false-lift
period theorem} replaces the repeated hyperbolic-plane estimates by
first occurrence. At each chosen finite place, two compact parabolic
orbits reduce all allowed rank-one data to the prescribed descent;
the single-row Fock type gives the real reduction. Thus the bottom
descent vanishes. A first nonzero lift to \(I_1\perp\mathbb H^j\)
would be cuspidal. Its fixed Schwartz seminorm gives one growth
exponent for every right derivative, so the compact-unipotent
Poincar\'e inequality proves rapid decay. In rank \(d=1+2j\),
the actual split Hecke equation is \(F_{n-d}(X_0)H_d(X')F=0\),
with \(F_{n-d}(X_0)\ne0\) and \(q^{d+1}|X'|\le1\).
The Gaussian-binomial identity gives
\(\|H_d(X')-1\|\le e^{1/2}-1<1\), contradicting first occurrence.

The full-rank Fourier coefficient of the resulting zero hyperbolic
descent is a weighted integral over orthonormal frames. Its exact
splitting character leaves no frame-dependent determinant factor.
Integral Witt transitivity, including at 2, makes the good-place
fibres invisible to the level-invariant period. Additive characters
separate the remaining compact frame image; positivity then forces
the desired block period to vanish. Absorbing the residual rank-one
theta kernel gives the false-lift test. The six-copy small block has
rank 14, so the required inequality is \(n=2z>2\cdot14+5\).

The short-form input (A) asks for a spanning family of integral
source vectors whose normalized images have squared norms at
most \(\exp(H-mR+E_A)\) at every relevant conjugate. Here \(H\)
is the main amount and \(R\) is the source's total remainder,
not yet the analytic remainder \(R_{q,d}\). The input also
controls the finite clearing index for every selected translate.
The baseline gauge and complete definite-chain norm estimate are
proved. Combining the exceptional corrections with the analytic
see-saw to obtain the normalized family in (A) remains open.

The following elementary part of input (B) is proved. Let
\(A_t:S\to T\) be maps between finite-dimensional Hilbert spaces
with probabilities \(p_t\), \(d_S=\dim S>0\), \(D=\dim T>0\).
Set
\[
 K=\sum_t p_t A_tA_t^*,\qquad
 \delta=d_S^{-1}\sum_t p_t\Tr(I-A_t^*A_t)_+.
\]
Suppose
\[
 \Tr K^2\le (1+\epsilon)d_S^2/D,\qquad
 e=2\delta+\epsilon<1,\quad \epsilon\ge0.
 \tag{4}
\]
\begin{proposition}[Spanning from two scalar moments]
For \(e<\eta<1\), at most
\[
 k=\left\lceil (2D/d_S)\log((1-e)/(\eta-e))\right\rceil
\]
maps span a subspace of dimension at least \((1-\eta)D\).
\end{proposition}
\begin{proof}
For a subspace \(F\subset T\) of dimension \(\alpha D\), let
\(U_t=A_t^{-1}F\). The positive source deficit bounds
\(\dim U_t-\|A_tP_{U_t}\|_{\rm HS}^2\).
Since \(A_tU_t\subset F\), averaging and using Cauchy--Schwarz
give
\[
 \sum_t p_t\dim U_t
 \le d_S\delta+\Tr(P_FK)
 \le d_S(\delta+\sqrt{\alpha(1+\epsilon)})
 \le d_S(1+\alpha+e)/2.
\]
Some map adds rank at least \(d_S(1-\alpha-e)/2\).
Also \(\delta d_S\ge d_S-D\), so \(d_S<2D\).
The missing share \(r=1-\alpha\) obeys
\(r_{\rm new}\le (1-\lambda)r+\lambda e\) with
\(\lambda=d_S/(2D)\), until it is at most \(e\).
Iteration from \(r=1\) proves the result.
\end{proof}

\paragraph{Equal good-form multiplicities.}
The \prooflink{notes/lsernau-good-form-multiplicities.md}{two-place
comparison} proves equality of the finite Hecke modules of good forms
on \(Y_0,W_1,\ldots,W_m\). Each good tuple has a primitive
cuspidal finite general-linear label at one place above each of the
two distinct level primes. A Cartan-intertwining calculation makes
both local components supercuspidal. The exact finite-type projector
is a sum of truncated supercuspidal coefficients. Harris--Labesse's
Proposition 1.10.6 and the refined argument of Theorems 3.1.5--3.1.6
therefore compare the specified real discrete series, without an
unverified stronger regularity hypothesis.

For signatures \((p_v,q_v)\), the real transfer sign is
\((-1)^{\sum_vp_vq_v}\). Finite isometry and the quadratic norm
product formula imply that \(\sum_v(q'_v-q_v)\) is even;
\(q(n-q)\equiv(n-1)q\pmod2\) makes the total relative sign one.
Normalized real minimal idempotents have rank one, so no compact-type
dimension or packet-cardinality factor intervenes. Varying the finite
tests, including their central cutoffs, identifies multiplicities of
the full finite representations, not just dimensions.

The \prooflink{notes/lsernau-common-moment-interface.md}
{finite dimension and moment argument} counts the definite
principal-level forms by class mass, using the triviality of
stabilizers at a split principal level. The proved comparison gives
the target bound \(D\le|\mathcal G|d_\infty\mathrm{mass}_r\).
For the source's scaled Poincar\'e map \(P\), the good
projector \(E\), and the finite group \(\mathcal G\), it
suffices to bound the unprojected moment
\(\sum_g\|P^*\rho(g)P\|_{\rm HS}^2\) and the positive
deficit of \(P^*EP\). Compression of the averaged positive
operator cannot increase its squared trace. The source's two numerical bounds, if established, give
\(e=19b^{-4}\) on the full source. The ramified source used
in (A) is the intersection of the nontrivial-character
parts of \(f\) commuting scalar actions at \(N\).
Free-level class counting gives its exact dimension
\[
 d_S=d_X(1-1/(N-1))^f.
\]
It is spanned by applying the integral operator
\(\prod_w(1-T_w)\) to the integral source generators.
The parameter inequality \(f/(N-1)\le b^{-8}\)
then permits restriction of both moments with
\(2\delta+\epsilon\le20b^{-4}\), as proved in the
common finite theta chapter. Thus at most
\(\lceil4D\log b/d_S\rceil\) translates of the actual
ramified source suffice for \(n\ge(1-4b^{-2})D\).
If the common local source colength estimates hold,
the selected translates have one-sided logarithmic
index at most \(5D\log(1+2b^2)fa(\log b)^3\).
The full see-saw and local colength budget remain obligations.
Absolute convergence and the local restriction estimate have the
following separate proof.

\paragraph{Quantitative restrictions and arithmetic returns.}
Write \(n_0=\rho+1\), \(r=k+2\rho+1\) and
\(f=[E^+:\mathbb Q]\) for the Bessel-frame parameters.
The \prooflink{notes/lsernau-quantitative-returns.md}{quantitative return
chapter} proves the Casimir filtration bound directly from compact
weights and unitarity. Schur orthogonality, the explicit Gram section
and a one-dimensional Hilbert-space estimate, iterated in the
\(2n_0^2\) real matrix coordinates, give
\[
 \int_{\|u\|\le P^{20}}
 \left\|\int_G\phi(y)W(\gamma_Z(uh)^{-1}y)\,dy\right\|^2
 d\mu_c(h,u)
 \le P^{214n_0^2}\frac{\|W\|_2^2\|\phi\|_2^2}{d_\pi}.
\]
Here \(W\) is reproduced by the genuine integrable discrete-series
coefficient, \(\phi\) has support \(\|y\|\le P^{20}\), and the
Gram ball, Haar measure, formal degree and parameter threshold are
exactly those specified in the chapter. Choosing the coordinate
width \(P^{-129}\) sharpens the source exponent \(2000n_0^2\).

For the stated integral frame, bounded reduction and free definite
class level, the product formula separates distinct frame images.
Their real dimension is \(2n_0rf\). Equal-frame terms form one
unipotent congruence coset, whose cutoff weights sum to at most one.
Packing and exact layer integration therefore give the uniform
weighted return bound
\[
 C_{\rm ret}\le
 2\bigl(1+4Nb^{74n_0}\bigr)^{2n_0rf}
 \le 2b^{(268n_0^2+2n_0)rf}.
\]
Combined with the Bessel majorant, this proves absolute Poincar\'e
convergence, the justified class-return Gram identity and the
truncation error \(\nu_U b^{-k^3f}/2\) in the stated source range.
The \prooflink{notes/lsernau-bessel-subgroup-moment.md}{exact Bessel-subgroup calculation} gives normalized subgroup matrices that are unitary on the reduced subgroup, zero outside it, and the identity at its identity element. Its summed squared Hilbert--Schmidt norm is the subgroup order times the source dimension. The full-minus-cut operator norm is at most $b^{-k^3 f}/2$. The non-Bessel second-moment estimates remain open.

\paragraph{Normalized lift and integral critical values.}
The \prooflink{notes/lsernau-normalized-transfer-critical-norm.md}
{normalized transfer chapter} supplies the full rational semisimple
representation identity for occurring comparison eigensystems,
the rectangular finite mass, and the stable-range norm formula
for critical pairs with $z\ge2$.
The \prooflink{notes/lsernau-integral-critical-kernels.md}
{integral kernel chapter} supplies their transported integral gauges
and group-ring measures, with denominator
$M^{fn+L}d^{\lceil L/2\rceil}$ in its rank-$n$, row-degree-$L$
notation. The two-place comparison now supplies occurrence of every
existing good eigensystem on the definite space. The rank-two endpoint
has a separate equal-rank period proof below; no stable-range theorem
is applied outside its hypotheses.

\paragraph{Good types directly in the Bessel model.}
The \prooflink{notes/lsernau-bessel-good-type-transfer.md}{finite
Bessel theorem} avoids the conjugacy-averaged local integral estimate
in this part of the proof. At a split place, the actual unipotent
group has blocks \((1^\rho,k+1,1^\rho)\) and one chain character,
with \(r=k+2\rho+1\). Andrews--Thiem's precise pattern-group
theorem changes it to the one-sided radical
\((1^{2\rho},k+1)\), with the same induced character and no scalar.
The unique generic representation of the mirabolic group in rank
\(a\) has dimension \(\prod_{j<a}(q_w^j-1)\): decompose along
its abelian first-row subgroup and induct on its nonzero character
orbit. Every generic \(GL_r\)-representation contains it.
Taking \(2\rho\) successive character spaces gives
\[
 \dim\operatorname{Hom}_U(\psi,\pi)
       \ge\prod_{j=1}^k(q_w^j-1),
\]
with equality for a cuspidal Green representation.

The \prooflink{notes/lsernau-finite-good-type-density.md}{finite
label counts} therefore apply directly to the probability weights
\(\mu_\psi(\pi)=|U|\dim\pi\dim\operatorname{Hom}_U(\psi,\pi)/|GL_r|\).
For residue size \(q_w\), rank \(r\ge2\), and forbidden root
exponent \(h_{\rm forb}\), put
\(a_q=1+\gcd(h_{\rm forb},q_w-1)\). The number \(A_r\) of
allowed squarefree monic labels has generating series
\[
 \sum_{r\ge0}A_r t^r
   =\frac{1-q_wt^2}{(1-q_wt)(1+t)^{a_q}}.
\]
Green's degree formula gives the sharp minimum Bessel label mass
\[
 \omega_B=\frac{\prod_{j=1}^k(1-q_w^{-j})}
                   {q_w^r(1-q_w^{-r})}.
\]
Exactly \(\varphi(q_w^r-1)/r\) labels are primitive, and all
are allowed. Their total mass is exactly
\(p_* =\varphi(q_w^r-1)\omega_B/r\).
For \(f\) factors the discarded share therefore obeys
\[
 \beta\le1-(A_r\omega_B)^f+(A_r\omega_B-p_*)^f
       \le f\varepsilon_*+e^{-fp_*},
 \quad
 \varepsilon_*=
 \frac{\gcd(h_{\rm forb},q_w-1)+2}{q_w}
                   +\sum_{j=1}^k q_w^{-j}.
\]
The already checked admissible inequalities give \(\beta_i\le2b^{-9}\)
at each chosen prime. A fixed class whose Bessel reduction is the
identity at the other prime has trivial definite component, by the
torsion-free principal congruence kernel. Thus the actual source trace
is \(d_X\psi\), and
\(\operatorname{tr}r_\psi(1-e_i)=d_X\beta_i\).

Write \(D_g=P^*\rho(g)P-r_\psi(g)\), with \(r_\psi\) extended by
zero outside the reduced Bessel subgroup \(R\). The remaining
restricted analytic estimate is
\[
 \sum_{g\in RG_i}\|D_g\|_{\rm HS}^2\le b^{-12}|R|d_X
                                      \quad(i=0,1).
\]
Equivariance and orthogonal convolution by a central idempotent
then bound each projected error by \(b^{-6}\sqrt{d_X}\).
Summing over \(R\) gives the same bound for \(P^*P-I\).
Taking traces proves
\[
 \operatorname{tr}(I-P^*EP)_+
       \le(3b^{-6}+4b^{-9})d_X\le4b^{-6}d_X.
\]
The identity and good-source deficits are therefore consequences of
one restricted-moment interface; the displayed analytic estimate
itself remains open.

The \prooflink{notes/lsernau-link-lattice-counting.md}{weighted lattice
theorem} closes its arithmetic counting input. Put
\(n_B=\rho+1\), \(\tau=b^{kf/4}\), \(H_c=b^{n_B^3f}\), and let
\(h_{\rm cls}\) count the free source classes. Write \(d(L)\) for
the normalized Hermitian discriminant. For the actual one-row
congruence lattices \(\Lambda_c\), it proves
\[
 \sum_c\sum_{\substack{L\subset\Lambda_c,\ \operatorname{rank}L=k-\ell\\
                       d(L)\le H_c}}d(L)^{-k}
 \le\frac{h_{\rm cls}}{N^{kf}}\tau^{\ell+1}.
\]
Minkowski gives a short basis of the saturated orthogonal complement;
Jacobi's identity preserves its discriminant. Packing bounds the
subspace height sum, and a local triangular normal form counts its
finite-index sublattices. Averaging over the \((N^k-1)^f\) nonzero
rows modulo the selected half of \(N\) gives the essential conductor
gain. The \prooflink{notes/lsernau-corank-frame-fibres.md}{frame-fibre
theorem} then uses three finite quotients: lifts modulo the Gram
lattice, the corank complement modulo \(M\delta\), and the isotropic
part modulo \(M\delta\bar\delta\). Explicit integral unipotent moves
identify each colour. The last column vanishes by its integral norm
bound \(4/N<1\). The count per lattice is
\(b^{2000n_B^6(\ell+1)f}\); squaring this fits one further factor
\(\tau\). Thus the full weighted near-pair count is proved for the
specified cell family, with exactly one class and one reduced-label
factor.

The \prooflink{notes/lsernau-link-gain-index.md}{link-gain proof}
supplies the congruence determinant gain and reciprocal torsion-index
cancellation; one precise integral link-lattice cut remains.
The \prooflink{notes/lsernau-single-link-energy-frontier.md}{single-link
reduction} proves the summation algebra. Given that cut,
the gauge-specific orbit estimate and reduction to
the specified nonzero cell family, the link energy over elements
with meet norm at least \(z\) is at most
\[
 (b^{2f}/z)^kD_{\rm main},\qquad
 D_{\rm main}=d_X^2/(d_\infty\mathrm{mass}_r).
\]
On either restricted strip, \(z\ge p_i^{2f}\ge(M/b)^f\).
The actual scale inequality
\(D_{\rm main}\le(M/b^6)^{kf}|R|d_X\), together with the proved
cutoff error, then gives \(9b^{-40}|R|d_X\), and source deficit
at most \(5b^{-9}d_X\). These sharper conclusions remain conditional
on the listed geometric suppliers; the proved pair count alone
does not close the restricted moments.

The \prooflink{notes/lsernau-positive-type-moment.md}{positive-type
moment chapter} controls the mass of singular leading minors by
counting their kernels with finite Fourier orthogonality.
It subtracts the guaranteed one-element kernel and retains the
geometric mean of the two stabilizer masses. Maximizing the weighted
fixer mass and using the source's open and terminal estimates then
gives the factor \(1+3b^{-4}\) instead of \(1+9b^{-4}\)
in its second moment. Those two actual unfolded estimates and the
restricted moments above remain required.

\section{Long forms from selected twists}
For input (C), the source forms \(\Theta_\xi\) must realize a
nonzero determinant \(Z(\xi)=\det(\Theta_\xi^*\Theta_\xi)\).
The quotient \(Z(1)/Z(\xi)\) compares Gram determinants on
\emph{the same basis}. One can choose \(\xi\) separately at
each finite place, but not apply the product formula to the
varying \(Z(\xi)\).

\subsection{The small-twist theorem}
The abstract theorem is supplied by the following
dependency-ordered proofs:
\begin{enumerate}
\item \prooflink{notes/lsernau-auxiliary-prime-supply.md}
{Auxiliary primes}: class field theory and Chebotarev construct
paired primes with the prescribed anticyclotomic characters.
\item \prooflink{notes/lsernau-classical-tied-descent.md}
{Tied Selmer descent}: ordinary Greenberg--Wiles duality and
reciprocity give short tied groups and a cardinality bound
independent of the later torsion level.
\item \prooflink{notes/lsernau-finite-selmer-presentations.md}
{Finite presentations}: a common polynomial matrix, sharp
Smith minors and a strict length cutoff give small Selmer witnesses.
\item \prooflink{notes/lsernau-twist-specialization.md}
{Torsion specialization}: a single tuple of roots of unity
works for arbitrarily many prescribed polynomials and signed powers.
\item \prooflink{notes/lsernau-general-small-twist.md}
{Combined theorem}: integral critical-value systems with Euler
factor norm relations have polynomial witnesses on a line.
All length excesses, value-valuation excesses, and tame inertia
costs over any finite set of signed powers can share one arbitrarily
small error.
\end{enumerate}
These are ordinary mathematical proofs under the explicit
Selmer datum and integral Stickelberger-system hypotheses.
They do not assert those hypotheses for the actual moving
automorphic representations or actual critical L-values.

\subsection{Attachment and the common row}
The \prooflink{notes/lsernau-galois-two-string.md}
{two-string attachment chapter} proves that two
norm-one torus characters can be chosen with their added
strings disjoint at every split finite place. Their weights
differ by the number of roots of unity at one complex place.
The separation follows from the ideal class group: a
norm-one idele with nonzero valuation at a split place
cannot have trivial value under their quotient.

For rank \(n\), the two lifts to the odd rank \(2n+1\),
if supplied with the source's attachment and local clauses,
have normalized traces \(\operatorname{tr} R_i\) differing
from the same row trace by the two disjoint character sums.
Linear independence of irreducible characters then proves
that the common difference is the trace of an actual
rank-\(n\) representation: for each irreducible class at
least one of the strings has multiplicity zero.
The row is a direct summand of both larger representations.
Unramifiedness, crystalline clauses and labelled jumps
pass to it, and the two coprime strings cancel in the
spherical annihilating polynomial. This supplies the
even-rank cancellation argument, conditional on the actual
theta lifts and odd-rank theorem. The
\prooflink{notes/lsernau-odd-rank-patching.md}{odd-rank patching
chapter} proves the extension from compatible representations on
open normal subgroups killed by a common set of inertia generators.
It reduces local labels to the split case using their actual
splitting degree, the least common multiple of the irreducible
factor degrees. The protected global steps and automorphic weak base
change have the suppliers described below. In the
common descent diagram, the exact normalized Hecke
polynomial identity, once established, also gives the
global semisimple decomposition of the lift into the
twisted row and its prescribed inverse-cyclotomic string.
The \prooflink{notes/lsernau-base-change-trace-transfer.md}{base-change trace}
proves the weighted groupoid and spectral transfer, and label extraction
from the matching-root, level, matching-product and finite Shintani inputs.
The \prooflink{notes/lsernau-specific-hecke-transfer.md}{specific Hecke transfer}
and \prooflink{notes/lsernau-shintani-label-transfer.md}{finite Shintani labels}
supply the matching products and labelled automorphic transfer using
the explicit classical finite-field character correspondence.
The \prooflink{notes/lsernau-global-norm-roots.md}{global norm-root chapter}
supplies the global matching roots and the matching compact families
by cyclic norm equations and finite free local polynomial algebras.

\paragraph{Spherical attachment at every rank and coefficient prime.}
The \prooflink{notes/lsernau-split-galois-attachment.md}{split attachment
chapter} verifies the classical input on the source Weil-data fields
at every positive rank and coefficient prime, including \(2\). The standard
definite unitary group is quasi-split
at every finite place: the CM quadratic extension is unramified
because the field contains \(i\) and \(\zeta_3\), and \(-1\) is
a local unit norm. Labesse's base change and discrete cuspidal
factors, followed by the explicit Tate and character twists of
Guerberoff's construction, reduce attachment to Chenevier--Harris's
any-prime cuspidal theorem. Labesse's compact-infinity hypothesis
holds, and the real degree is divisible by four since the field contains
\(\mathbb Q(\zeta_{16})^+\). The given finite-unramified Weil splitting
character supplies the parity characters in the cuspidal construction.
Thus no additional imaginary-quadratic splitting or odd-rank
condition is needed here. Local unramified base change and the
monodromy bound force actual unramifiedness at any spherical finite
place away from the coefficient prime; the crystalline clause gives
crystallinity above it. At split places the exact Frobenius polynomial
is the one displayed below.

The conjugate Weyl coefficient, geometric reciprocity and Hodge
sign are checked explicitly. For highest weight
\(\lambda_1\ge\cdots\ge\lambda_n\), the two labelled jump sets are
\(\{j-1-\lambda_j\}\) and \(\{\lambda_j+n-j\}\).
The source's unnormalized spherical operators give roots
\(q^{(n-1)/2}\alpha_i\), exactly those of the parameter twisted by
\(\lvert\det\rvert^{(1-n)/2}\). Cayley--Hamilton gives the required
polynomial on the actual matrix or crystalline period endomorphism.

For an unramified sorted Iwahori eigencharacter, finite column-orbit
counting and a minimal-index argument prove that spherical averaging
is nonzero; a hyperplane count gives its exact Hecke polynomial.
The tame descent index is
\(k\lvert\operatorname{im}(\mathcal O_F^\times\to(\mathbb C^\times)^n)\rvert
\le k(q-1)\), and each prime cyclic step divides it by that prime.
Finally, a noncommutative exponential-polynomial recurrence shows
that the positive Frobenius-power identities imply the inertia
annihilator; at most
\((n(n-1)+1)\binom{2n-1}{n}\) powers suffice.
The \prooflink{notes/lsernau-iwahori-dominant-matching.md}{dominant
matching proof} twisted-conjugates each element of the strict Iwahori
double coset to a diagonal element by nested compact congruence conditions. Polynomial
interpolation then gives one matching root for every matrix with the
specified characteristic polynomial; principal-unit norms and the
trace-one Hilbert--90 argument preserve the exact support and labels.
The \prooflink{notes/lsernau-iwahori-type-sorting.md}{finite-type
extraction} starts from the prescribed parabolic product character.
A rank-two averaging identity proves injectivity of the dominant
operators and exchanges every out-of-order pair of eigencharacters.
Only strict inversions are exchanged, so ties and repeated labels
cause no division by a vanishing root difference.
The \prooflink{notes/lsernau-shintani-product-types.md}{full finite
Shintani product} preserves the determinant-block multiplicities
through genuine semidirect parabolic induction.

The \prooflink{notes/lsernau-tame-local-global-assembly.md}{prime-tower
assembly} supplies the global steps and descent. Prime-exponent
character approximation puts auxiliary ramification at places of
rational-prime norm, hence at CM-split places. Symmetrizing the
character gives a degree-\(p\) CM extension, also for \(p=2\);
it is unramified above \(p\), has an integral trace-one element,
and admits the required extension of the Weil data by a half-different
supported on conjugate pairs. The auxiliary character is trivial
on the whole distinguished decomposition group.

For each chosen positive Weil element \(s\), the extension is
chosen to keep \(s\) in
\(H=\ker(\nu+\nu^c)\cap\ker\nu\).
Weak base change and Chebotarev identify the two semisimple
Galois representations on \(H\). The intertwiner transports the
actual polynomial operator, including on potentially crystalline
periods. The induction reaches the spherical case in a tower of
total degree exactly \(k|C|\le k(q-1)\).
The finite all-powers lemma gives the inertia identities.
A tower unramified at the distinguished place first splits a
squarefree Green label; the
product-type transfer, sorted extraction and same descent then
give its full labelled annihilator with the prescribed repeated
determinant factors. Finite inertia away from the coefficient prime,
potential crystallinity above it, and the labelled jumps descend
at the same time.

\subsection{From the attached main summands to Selmer data}
The \prooflink{notes/lsernau-main-summand-data.md}{datum construction}
starts with attached representations having the source's local
type and Hodge compatibility. The primitive labels yield a tame
inertia operator with a full Frobenius orbit of eigenlines, distinct
even modulo the coefficient prime. Integral eigenprojectors and
Frobenius then supply all matrix units, as well as a scalar at a
prescribed cyclotomic value. A large-order scalar inertia element
normalizes the cocycles. The conjugate-dual defect is at most
\(v_\ell(G)\), by filtering its intertwining matrices and applying
residual Schur's lemma in each graded piece. Local Tate duality
bounds the local terms, and the crystalline rank formula bounds
the remaining local conditions. For a main summand of rank \(n\),
this gives the normalized datum defect
\[
 B_i\le ((4f+2)n+1)v_\ell(G),
\]
where \(G\) is the chosen finite group order. The source's bound
\((4f+3)r v_\ell(G)\) follows. This argument constructs the
Selmer data \emph{from} the attached type-compatible representations.
Attachment, the normalized rational theta-lift identity and occurrence
on the definite comparison space are supplied above. Integral lattice
splitting and nonvanishing of the full chain do not follow from that
rational identity.

\subsection{Critical values and the chain correction}
The \prooflink{notes/lsernau-critical-value-systems.md}
{integral-system construction} expands the squared norms of
the rank-\(2z\) character lifts for \(z\ge2\), and the
rank-two determinant period square for \(z=1\), into group-ring
elements. The finite stabilizer denominators divide \(G\);
an idelic Hilbert 90 pushforward produces the exact
Stickelberger systems from these measures, with no division
by a torus class number. The integral-kernel chapter supplies the
stable-range gauge and norm formula for \(z\ge2\).

\begin{samepage}
For \(z=1\), the
\prooflink{notes/lsernau-rank-two-period-arithmetic.md}{rank-two
arithmetic theorem} observes that the tame determinant datum factors
through the radical of its stated modulus. If
\(h=2\sum_\sigma(c_\sigma-1)\), its uniform integral scale is
\[
 S_\nu=\mathrm{vacuumScale}^{4}\mathrm{jetScale}^{h}
                       M^{4f+h}d^{h/2}.
\]
\end{samepage}
The identity-residue Fourier expansion has exactly \(4\) matrix
coordinates over the chosen half modulus. The charged-jet bound and
rational unitary invariance then give integrality at every coefficient
embedding away from the auxiliary norm \(N_Q\). Clearing stabilizers
by \(G\) gives an integral measure with evaluations
\(A_Q(\tau)=GS_\nu\Lambda_Q(\tau)\); multiplying by its conjugate
inverse gives \(F_Q(\tau)=|GS_\nu\Lambda_Q(\tau)|^2\).
The scale saves the integer factor
\(M^{12f+h}d^{h/2}\) from the older sufficient scale.

The finite \(q\)-binomial identity inserts the invertible-matrix
indicator. For an added unramified place of residue order \(q_w\),
put \(y=(\nu\tau)(\varpi_w/\bar\varpi_w)\). Invariance of the
quotient integral, without assuming that individual translates
preserve the old level, gives exactly
\[
 \Lambda_{Q\cup\{w\}}(\tau)
   =(1-y^{-1})(1-q_w^{-1}y^{-1})\Lambda_Q(\tau).
\]
Its absolute square is the required Euler norm relation. It can
vanish for imprimitive twists with \(y=1\).
The \prooflink{notes/lsernau-rank-two-period-norm.md}{rank-two analytic
period theorem} now proves, for every fully ramified tame twist,
\[
 |\Lambda_Q(\tau)|^2=C_\nu N_Q
       |L^{M\cup Q}(1,(\nu\tau)^{-\mathrm{BC}})|^2,
 \qquad C_\nu>0,
\]
with \(C_\nu\) independent of \(Q,\tau\), and proves the nonzero
initial period. Its proof has three steps.

First, equal-rank Hecke transfer makes the theta lift of the constant
function a \((q+1)\)-eigenvector for an operator that is a sum of
\(q+1\) unitary translations. Equality in the norm bound makes each
translation fix the lift. Strong approximation for \(SU(2)\), together
with weak approximation for the determinant torus, then makes the lift
constant. Exact transposition of the positive Hom space identifies its
value with the prescribed period; no swap phase remains.

Second, anisotropic Siegel--Weil applies at \(m=n=2\), \(s=0\):
its multiplier is 2, and no regularization is needed. Every rational
maximal isotropic subspace of the doubled space is the graph of an
isometry, so doubling unfolds through one orbit. The determinant
correction cancels on the diagonal. For
\(\chi=(\nu\tau)^{-\mathrm{BC}}\), direct split and inert Cartan sums
give the continued finite product
\[
 \frac{L_{E,f}(0,\chi)L_{E,f}(1,\chi)}
      {L_{F,f}(1,\eta)\zeta_{F,f}(2)}.
\]
At a ramified split place the matrix coefficient is supported exactly
on the compact group and its integral is
\((1-q^{-1})(1-q^{-2})\), the same denominator at zero.
The real determinant vector has positive norm \(k!(k+1)!\).
All Haar conversion factors are fixed and positive.

Finally the Hecke functional equation, with infinity shift
\(c_\sigma\ge1\) and square-root conductor
\(|D_E|^{1/2}M^fN_Q\), gives the displayed period square.
The Euler-product inequality from
\(3+4\cos\theta+\cos2\theta=2(1+\cos\theta)^2\)
proves \(L(1,\chi)\ne0\): the fourth power of a hypothetical zero
cannot be canceled by the third power of the zeta pole, while
\(L(s,\chi^2)\) is holomorphic at 1. Thus the endpoint construction
and its nonzero base case are supplied without three complementary
theta exchanges.

The initial tame value \(\lambda\ne1\) is a root of unity.
Its exact valuation \(v_\ell(1-\lambda)\) is zero unless its
order is an \(\ell\)-power, when it is at most
\(1/(\ell-1)\). In the small-twist theorem the base measure
evaluation cancels, and a critical ratio \(R_\xi\) satisfies
\[
 v_\ell(\iota R_\xi)\le2v_\ell(1-\lambda)+2^{-\ell}.
\]
Choosing a fixed initial place above \(p_0\) except at
\(\ell=p_0\), when a place above \(p_1\) is used, gives
a total inverse-ratio height less than
\(4\log 2+3/2<5\) per critical pair, in the source's
height convention. This is conditional on the actual
measure inputs just stated. It is uniform over the
coefficient isomorphisms and the auxiliary level.

With the common chain and the single-lift Rallis Gram
identities, determinant multiplicativity gives
\[
 \frac{Z(\xi)}{Z(1)}
 =\left(N_Q^{e/2}\prod_{i<j}\mathcal R_{ij}(\xi)\right)^D
       \prod_{j=1}^m\frac{\Delta_j(\xi)}{\Delta_j(1)},
\]
where \(\Delta_j\) are the single-lift volumes, the first
factor is the exchange correction, and every volume must be
nonzero. The correction before exponent \(D\) has inverse
height at most \(5t\,m(m-1)/2\) under the above hypotheses.
The stable-range single-lift Rallis identity is supplied above under
its stated real-edge hypotheses. The full-chain exchange identity,
its compatible nonzero volumes and its quantitative application
remain to be established.

\paragraph{The algebraic congruence element.}
The \prooflink{notes/lsernau-ribet-polynomial.md}{Ribet polynomial
chapter} proves the common tied-row element by an integral
determinant expansion over primitive cyclic words and
Cayley--Hamilton. It gives an exact product of tied-row block
determinants at each lift, with no new factorial denominators.
The \prooflink{notes/lsernau-tied-row-construction.md}{tied-row
construction} gives the main and character matrices by DVR
duality and elementary divisors, with the exact common squarefree
coefficient and its valuation budget. The arithmetic length input
has a direct reduction to local crystalline defects, as follows.
For an integral lattice \(T\), let \(Z_f(T)\) be the cocycles whose
rational classes lie in \(H_f^1\), including coboundaries. Put
\(T^*=\operatorname{Hom}(T,O)(1)\). Local Tate duality, the Euler
characteristic and the complementary Hodge weights give
\[
 |\overline{Z_f(T)}|\,|\overline{Z_f(T^*)}|
       =|T/p^aT|^{[E_v:\mathbb Q_p]+2}.
\]
The \prooflink{notes/lsernau-integral-separators.md}{integral separator
chapter} includes the coboundary factors in the finite duality count.
A Smith-basis construction lifts every relation-respecting Ribet
cocycle through a full lattice in the same prescribed local class.
Its modular cup products are therefore controlled by the dual defect.
Restriction and the normalized tame-generator conditions yield
\[
 \ell_p(Z_{\rm Ribet})\le
 \ell_p(\mathrm{Sel})+
 \sum_{v\mid p}\ell_p D_{\mathcal Q_v}(T^*),
 \qquad \ell_p=\operatorname{length}_O/v_O(p).
\]
This proves the passage from local defects to tied-row lengths,
including divisible coefficients. The numerical bound on these
defects remains a distinct arithmetic input.

\paragraph{What qualitative crystalline control supplies.}
For a fixed potentially crystalline family \(u_y\), put
\(A(g)=(u_y(g)-1)_y\), \(I=\sum_g OA(g)\), and
\(H=\operatorname{Hom}_O(I,T)\). The possibly unsaturated lattice
\(I\) is used as it stands. Its universal cocycle extension embeds
rationally in \(\bigoplus_y\operatorname{Mat}_{n_y}(F)\oplus F\),
so crystalline exactness and Galois descent make \(A\) crystalline.
Every relation-respecting modular cocycle is \(q\circ A\).
The \prooflink{notes/lsernau-fixed-relation-smith-bounds.md}{Smith theorem}
therefore bounds its crystalline obstruction by
\[
 \kappa_{T,v}=e_O^{-1}\sum_i a_i,
\]
where \(a_i\) are the nonzero Smith valuations of an invariance matrix
on \(H\). One maximal-rank minor with certified rational kernel
already gives an upper bound. The
\prooflink{notes/lsernau-uniform-inertia-relations.md}{inertia-algebra
version} removes dependence on multiplicity and auxiliary precision
when the rational inertia algebra is fixed.

The stronger
\prooflink{notes/lsernau-bounded-class-crystalline-loss.md}{bounded-class
theorem} fixes only the base field, target lattice, crystallizing
extension and weight interval. In the Artinian quotient
\[
 H^1(K,V/T)\big/\operatorname{im}H_f^1(K,V),
\]
the subgroup represented by torsion-crystalline extensions has finite
length. Otherwise its divisible part gives a coherent obstructed tower;
finiteness at each precision lifts that tower to one integral cocycle.
Liu's torsion-limit theorem makes its rational extension crystalline,
a contradiction. Modular pushouts of the universal relation extension
belong to this class. Thus the bound works for the entire enlarged
family \(Y\), at every rank and precision, including 2 and level primes;
coefficient extension preserves normalized length. This supplies
finiteness for fixed data, without a numerical bound as the data grow.

\paragraph{An exact loss and the importance of sparse weights.}
The \prooflink{notes/lsernau-exact-cyclotomic-loss-and-budget.md}
{cyclotomic calculation} shows why the last distinction is necessary.
Over \(\Q_p\), take trivial target and auxiliary characters
\(\chi^{d_pj}\), \(1\le j\le N\), where
\((d_p,a_p)=(p-1,1)\) for odd \(p\) and \((2,3)\) for \(p=2\).
Choose inertia with \(\chi^{d_p}(\gamma)=q=1+p^{a_p}\).
The actual relation lattice is cyclic for
\(U=\operatorname{diag}(q,\ldots,q^N)\), with polynomial
\(P_N(X)=\prod_j(X-q^j)\). Its coinvariants are
\(\Z_p/P_N(1)\). Evaluation at \(\gamma\) identifies the
noncrystalline classes: integral crystalline cocycles for the trivial
target are unramified, whereas these are normalized ramified logarithms.
Consequently their exact length modulo \(p^n\) is
\[
 \min\{n,a_pN+v_p(N!)\}.
\]
Even with fields and target fixed, a universal bound using only interval
width must grow at least linearly. This is not a counterexample to the
source's sparse labelled family; it shows why that information matters.

For the restricted class whose rational inertia action is a sum of
\(1\) and \(\chi^s\), \(s\in S\subset\Z\setminus\{0\}\), a
single annihilating polynomial gives the explicit all-prime bound
\[
 \sum_p\kappa_p\log p\le\log Q(S),\qquad
 Q(S)=2^{1+2|S|}\prod_{s\in S}
       \left(|s|\prod_{d\mid |s|}(d+1)\right).
\]
For \(S=\varnothing\), take \(Q(S)=1\). This bound is independent
of multiplicities, lattices and precision and fits the source scales
under the stated sparse-size inequalities. Crystalline weights alone
do not imply this cyclotomic inertia decomposition for the actual
nonabelian family. Its effective cost remains an arithmetic obligation.

\paragraph{Zero defect in the Fontaine--Laffaille range.}
The \prooflink{notes/lsernau-fontaine-laffaille-zero-defect.md}{zero-defect
chapter} supplies a complete large-prime subcase. Suppose the base
field \(E_v/\mathbb Q_p\) is unramified, the class is crystalline,
and at each labelled embedding its allowed weights together with
zero lie in an interval of length at most \(p-2\).
A labelled rank-one crystalline twist moves the two lattices and
the extension's rank-one quotient into one Fontaine--Laffaille range.
Torsion full faithfulness then transports every map modulo \(p^a\).
The resulting extension lifts with both integral endpoints fixed:
in adapted filtered bases its normalized Frobenius matrix is
\[
 \begin{pmatrix}A_M&B\\0&A_N\end{pmatrix},
\]
where \(A_M,A_N\) are the given invertible endpoint matrices and
only the off-diagonal block \(B\) must be lifted. Strong divisibility
is preserved. Therefore every modular map preserves the integral
finite Bloch--Kato classes, and local orthogonality gives
\(D_{\mathcal Q_v}(T^*)=0\).

For the source's cyclotomic field and aligned weight windows, this
applies whenever \(p\nmid M\) and
\(p\ge b+\delta+W+\eta+2\). In particular it covers the entire
zero-loss branch \(p\nmid M\), \(\lambda_p=0\), because
\(\lambda_p=0\) implies \(p>2b\), and all conductor primes are
at most \(Q<b\). Explicit level-prime and positive-loss bounds on
the growing global data are still required. At the primes covered here the tied length pays
no crystalline defect term.
The \prooflink{notes/lsernau-local-tied-derivatives.md}{local derivative}
and \prooflink{notes/lsernau-tied-integral-extraction.md}{integral extraction}
proofs convert labelled attachment into a denominator-clearing operator.
\prooflink{notes/lsernau-triangle-detection.md}{Triangle detection}
uses at most \(y^4\) Frobenius elements;
\prooflink{notes/lsernau-kummer-relations.md}{Kummer duality} and
\prooflink{notes/lsernau-root-chain-selectors.md}{root-chain selectors}
provide its strict dimension and separation inputs, including at 2.
The numerical defect, Selmer, triangle and precision bounds remain
hypotheses. For a fixed family the Smith theorem gives
\[
 B_p^{\rm fixed}=s+2t'a+\sum_i\sum_{v\mid p}\kappa_{T_i,v}
                     +\operatorname{denominatorError}_p.
\]
Each actual main or character summand is charged once. The bounded-class
theorem permits the replacement
\(\sum\kappa_{T_i,v}\mapsto e_O^{-1}\sum\operatorname{length}_O C_{i,v}\),
which is fixed before enlarging the auxiliary family or precision.
This removes the family-versus-precision circularity of the raw Smith
cost. All other tied-row and extraction inequalities remain, and the
new constants have not been bounded on the growing global data.
Over an
unramified coefficient extension with residue field larger than the
number of eigensystems, one linear combination preserves the best
valuation in every column: avoid the corresponding finitely many
proper hyperplanes. Thus simultaneous assembly has no additional
valuation loss. A stricter spectral projector has instead the exact
cost given by the largest Smith exponent of its congruence module;
qualitative semisimplicity alone does not bound that cost.

\subsection{The one-lift comparison}
For the exact descent pair at step \(j\), let \(L_\xi\) be the
lift on the good forms, \(G_\xi\) the finite gluing map of the
first factors and second-factor tuples, and \(h\) the source
Fock norm of the highest polynomial in the second-factor weight.
The \prooflink{notes/lsernau-one-lift-comparison.md}
{one-lift chapter} constructs its integral tight frame
using the Fock pairing and the column-filling matrix.
The comparison operator is
\(U_\xi=g_0 h^2G_\xi G_\xi^*\), with \(g_0\) the integer
clearing class stabilizers. In the finite gluing sum the
multiplier \(M^{l_1}\) cancels the divided-splitting gauge
cost \(|M|_\ell^{-l_1}\). At infinity Cauchy--Schwarz
on tuple entries bounds the gluing norm; in rational value
metrics, the adjoint at a conjugate is the adjoint of the
map at the corresponding conjugated automorphism. These
finite estimates leave the actual first-factor and local
product-form estimates as inputs. Write \(C_\xi\) for its compression
to \(Z_\xi=\operatorname{im} L_\xi\), \(\Delta_\xi=\det C_\xi\).
The required period identity
\(c_\xi G_\xi^*L_\xi=P^\circ\), with \(P^\circ\) fixed
and injective, gives the exact squared-volume comparison
\[
 \frac{\det(L_1^*L_1)}{\det(L_\xi^*L_\xi)}
   =\frac{\Delta_\xi}{\Delta_1}
     (\Gamma_1/\Gamma_\xi)^d,
\]
where \(\Gamma_\xi\) is the source three-exchange Hecke
product and \(|c_1/c_\xi|^2=\Gamma_\xi/\Gamma_1\).

The preceding first-order congruence construction supplies a
diagonal clearing map for each lift generator once its numerical
arithmetic and extraction hypotheses hold. Their uniform verification
must give the required valuation bounds \(b_i\). An integral linear combination
avoiding finitely many residue-field hyperplanes gives one
simultaneous operator \(\Lambda\) with
\(v_\ell(\det\Lambda)\le\sum_i b_i\).
The stated integral product-form gluing inputs and the still-required
higher-rank false-lift periods for the general comparison make
\(\Lambda C_\xi\) preserve the integral true lifts. The proved
rank-one period supplies the sharp critical test above; it does not
replace these higher-rank periods. Scaling eigenvectors to primitive integral
coordinates shows \(v_\ell(\Delta_\xi)\ge-\sum_i b_i\).
The false-lift test is also what ensures that the original
orthogonal projection contracts the norm at every complex
automorphism.

Let \(B\) be the denominator budget and \(J-P\) the source
direction amount (jet size less prime gain).
The signed main critical-ratio estimate follows by the
rank-one projected pairing and the actual sharp test-form
estimate. With the constant inverse-ratio bound of less
than 5 for its auxiliary rank-four pair, it costs at most
\(4B+10-(J-P)\), replacing \(6B-(J-P)\) when \(B>5\).
The three-exchange product then costs
\(50m+k(4B+10-(J-P))\), retaining the negative main term.
The full derivation, the exact test conditions and the
remaining analytic, integral and congruence inputs are in
the one-lift chapter.

\paragraph{Only one aggregate projection budget is required.}
There is no need to impose the same worst-case bound on every
projection. Put \(L=\log b\), \(B_0=faL^3/(mk)\),
\(A=3fmbL^2\), and \(\Delta=J-P\). Let \(\Lambda_j\) be the
all-prime log determinant cost of the full projection at lift \(j\),
and \(C_{j,\alpha}\) the height of its projected pairing at main
critical pair \(\alpha\); there are \(k/2\) such pairs. Define
\[
 \mathcal A=d^{-1}\sum_{j=1}^m\Lambda_j
       +2\sum_{j=1}^m\sum_{\alpha}C_{j,\alpha}.
\]
The \prooflink{notes/lsernau-exact-cyclotomic-loss-and-budget.md}
{aggregate budget identity} sums the actual four-factor comparison.
If the norm and each sharp-test cost are at most \(B_0\), the inverse
chain determinant has height at most
\[
 d\{B_{\rm chain}+mA+50m^2+mk(3B_0+10)
                      +\mathcal A-mk\Delta\}.
\]
Thus it suffices to prove
\[
 \mathcal A\le faL^4-mA-50m^2-mk(3B_0+10).
\]
All Selmer, detection and local denominator costs must enter this sum.
The negative direction term is preserved. Qualitative finiteness, or
increasing auxiliary precision, does not imply this inequality.

\section{The final comparison}
Here is the proved scalar argument. Suppose an independent
family \(y_1,\ldots,y_n\) in the target is selected, \(0<n\le D\).
Complete it to a basis. Write \(g\ne0\) for its full untwisted
Gram determinant. In the rational terminal value metric \(H_0\),
if \(W\) is its matrix of columns and \(c\) is complex
conjugation, then
\[
 \sigma(g)=\det((c\sigma c W)^*H_0(\sigma W)).
\]
Exterior Cauchy--Schwarz bounds this by the geometric mean
of the two squared covolumes. Uniform bounds at all complex
automorphisms therefore make \(g\) algebraic: a transcendental
number can be sent to arbitrarily large values by extending a
transcendence basis. The common height multiplier is algebraic
for the same reason. Choose a number field \(F\) containing
these common numbers; \(C\) is a bound
for the selected logarithmic squared covolume, and \(Q_0\)
bounds the complementary archimedean cost. At each finite
place \(v\) suppose \(g=u_v r_v\) with both factors nonzero.
Local integrality bounds \(\log|u_v|_v\le i_v\).
Suppose a \emph{single} \(h\in F^\times\) and nonnegative
\(b_v\) satisfy \(\log|hr_v|_v\le b_v\), with
\[
 I=\sum_{v\ {\rm finite}}\omega_v i_v,\qquad
 B_0=[F:\Q]^{-1}\sum_\sigma\log|\sigma h|
       +\sum_{v\ {\rm finite}}\omega_v b_v,
 \quad \omega_v=[F_v:\Q_p]/[F:\Q].
\]
Here absolute values extend the usual normalized ones, and
the sums include all places. Apply the product formula to
\(g\) and \(h\). The result is
\[
 0\le C+I+Q_0+B_0.
 \tag{5}
\]
In the actual Gram comparison clearing both conjugate sides
gives \(I=2\delta_{\rm clear}\). The correction \(r_v\)
includes the determinant ratio and the complementary clearing
factors; these must be combined before bounding their heights.
See the \prooflink{notes/lsernau-mathematical-extraction-2026-11-26.md}
{determinant proof} for the basis and clearing identities.

Hadamard and the lower determinant supplier require
\(C\le n(H-mR+E_A)\), \(B_0\le-DH+E_B\).
Substitution into (5) gives the exact signed test
\[
 nmR\le -(D-n)H+nE_A+E_B+I+Q_0.
 \tag{6}
\]
The missing-rank term is favorable only when \(H\ge0\);
it cannot be discarded otherwise.

For orientation, the readable source provides the sufficient
interface \(H\ge0\), \(n\ge(1-4/b^2)D\), and
\[
 \begin{split}
 C&\le n(H-mR+fa\log^4 b),\\
 D(H-3fa\log^4 b)&\le C+2\delta_{\rm clear},\\
 \delta_{\rm clear}&\le 10D\log(1+2b^2)fa\log^3 b.
 \end{split}
 \tag{7}
\]
On its scale range this yields \(mR\le fa\log^5 b\):
for sufficiently large \(b\), \(n\ge D/2\);
after division by \(n\) the sum of errors is
\(O(fa\log^4 b)\le fa\log^5 b\).
The inequalities (7) are the source's arithmetic suppliers,
\emph{not} new proven estimates for the objects in Section 3.

\section{From the arithmetic comparison to the theorem}
The exact source prime identity (proved in the
\prooflink{notes/lsernau-moment-and-prime-reductions.md}
{fixed-modulus reduction}) gives from (7)
\[
 R_{q,d}(x)\le \frac{\varphi(q)(\log b)^6}{m}
     +|R_{q,d}(x_1)|+\log b\,|R_{1,d}(x_1)|.
 \tag{8}
\]
The classical prime number theorem, with \(q,d\) fixed and
\(\log x_1\ge x^{\varepsilon/12}\), makes the last two terms
smaller than every power of \(x^{-1}\). Since \(m\ge\sqrt x\)
and \(\log b\ll x^\varepsilon\), a smaller accuracy parameter
gives (1); integer-to-real interpolation costs \(O(1/x)\).
The remaining work is the compatible arithmetic input.

